\section{Introduction}
\label{sec:intro}

In the era of large-scale foundation models, the machine learning community has witnessed a remarkable proliferation of specialized, task-specific expert networks. These models are typically fine-tuned from a shared pre-trained base model to excel at distinct, isolated domains (e.g., classification, segmentation, or specific language tasks)~\cite{ilharco2022editing, wortsman2022model}. While deploying separate expert networks for each task is computationally viable, it scales poorly in multi-task, resource-constrained environments due to the multiplicative growth of storage and inference memory footprints. 

To resolve this bottleneck, model merging (or parameter-space model fusion) has emerged as a compelling, training-free paradigm~\cite{yadav2023ties, matena2021merging}. By algebraically combining the parameters of task-specific experts, model merging aims to construct a single, unified multi-task network that retains the specialized capabilities of all constituent models simultaneously, without requiring expensive retraining or fine-tuning.

Early work in this domain relied on static merging rules, such as simple uniform averaging or heuristic coordinate alignment (e.g., Task Arithmetic~\cite{ilharco2022editing} and TIES-merging~\cite{yadav2023ties}). However, these static techniques suffer from severe performance degradation when combined models exhibit high parameter-space conflicts. To address this, active model merging methods, such as AdaMerging~\cite{adamerging}, optimize the merging coefficients dynamically. In online test-time adaptation (TTA) scenarios, AdaMerging minimizes the prediction entropy on unlabeled test-time data streams to adapt the layer-wise model weights on-the-fly.

Despite its conceptual appeal, unconstrained online AdaMerging is plagued by the \emph{Overfitting-Optimizer Paradox}~\cite{regcalmerge, polymerge}. Because the adaptation optimization is performed over a tiny, local, and unlabeled calibration set (or incoming test stream) of perhaps only a few dozen samples, optimizing high-dimensional layer-wise coefficients (e.g., $L \times K$ parameters for $L$ layers and $K$ tasks) easily results in transductive overfitting. Rather than finding a generalizable balance between experts, the optimizer fits the local noise of the unlabeled stream, causing a catastrophic representational collapse across all tasks.

To combat this overfitting, recent literature has introduced increasingly Byzantine and convoluted regularization schemes. For instance, RegCalMerge~\cite{regcalmerge} introduces Class-Capacity Normalization (CCN), Scale-Normalized Entropy Weighting (SNEW), and Elastic Spatial Regularization (ESR) with multiple hyperparameters. PolyMerge~\cite{polymerge} parameterizes the merging trajectory as a continuous, low-degree polynomial subspace over the network depth. Other concurrent works like QWS-Merge introduce quantum wave analogies using frozen random projections, normalized phase states, and interference equations to smooth the optimization landscape.

Guided by the philosophical principle of Occam's razor, we argue that this rapid escalation of complexity is entirely unnecessary. In this work, we demonstrate that the Overfitting-Optimizer Paradox is not an inherent flaw that requires heavy-handed spatial regularizers or continuous geometric parameterizations to solve. Instead, it is a straightforward consequence of excessive optimization degrees of freedom on low-sample regimes. 

We propose a simple, training-free, and highly elegant alternative: \textbf{Pruned Gradient Merging (PG-Merge)}. PG-Merge resolves test-time model fusion overfitting by applying a dynamic, non-parametric sparse gradient mask during adaptation. Rather than updating all merging coefficients, PG-Merge flattens the raw coefficient gradients, sorts their absolute magnitudes, and zeroes out all but the top-$p\%$ (e.g., top-$15\%$) most critical values. Standard optimizer updates (e.g., Adam or SGD) are then applied only to these highly sensitive coordinates, while the remaining $(100-p)\%$ of the parameters are kept strictly frozen.

The core contributions of this work are as follows:
\begin{itemize}
    \item We expose the redundant complexity in current test-time model merging regularizers and show that their performance can be matched or exceeded through a minimalist, non-parametric gradient-sparsity mechanism.
    \item We introduce PG-Merge, a training-free and hyperparameter-lean framework that dynamically filters out transductive noise and limits optimization degrees of freedom to prevent overfitting.
    \item We perform exhaustive evaluations on a compact Vision Transformer (\texttt{vit\_tiny}) across four diverse vision tasks (MNIST, FashionMNIST, CIFAR-10, and SVHN), showing that PG-Merge substantially outperforms unconstrained AdaMerging and matches or exceeds the generalizability of state-of-the-art regularizers.
    \item We provide a detailed ablation study characterizing the "sparsity sweet spot" for test-time adaptation, confirming that limiting active updates to $15\%$--$30\%$ of coefficients represents an optimal low-pass filter for parameter-space model fusion.
\end{itemize}